\documentclass[aspectratio=169, 14pt]{beamer}

\usepackage[T2A]{fontenc}
\usepackage[utf8]{inputenc}
\usepackage[russian]{babel}
\usepackage{tempora}
\usepackage{graphicx}
\usepackage{tikz}
\usetikzlibrary{shapes.geometric, arrows.meta, positioning}
\usepackage{listings}
\usepackage{tabularx}
\usepackage{booktabs}
\usepackage{array}

\usetheme{Madrid}
\usecolortheme{whale}

\setbeamertemplate{navigation symbols}{}
\setbeamertemplate{footline}[frame number]

\lstset{
  basicstyle=\footnotesize\ttfamily,
  breaklines=true,
  frame=none,
  numbers=none,
  tabsize=2,
  showstringspaces=false,
}

\title{Разработка RESTful веб-приложения\\для управления задачами}
\subtitle{на основе многослойной архитектуры\\с использованием языка программирования Go}
\author{Берсиров~С.}
\institute{ДГТУ, Факультет информатики и ВТ\\Кафедра информатики}
\date{2026}

\begin{document}

% ============================================================
% Слайд 1 — Титульный
% ============================================================
\begin{frame}
\titlepage
\end{frame}

% ============================================================
% Слайд 2 — Цель и задачи
% ============================================================
\begin{frame}{Цель и задачи}

\textbf{Цель:} разработка RESTful веб-приложения TodoApp для управления задачами на языке Go с многослойной архитектурой.

\vspace{0.5cm}
\textbf{Задачи:}
\begin{itemize}
  \item Провести анализ предметной области и аналогов
  \item Обосновать выбор технологий и архитектуры
  \item Спроектировать БД и REST~API
  \item Реализовать серверную часть приложения
  \item Провести функциональное тестирование
\end{itemize}
\end{frame}

% ============================================================
% Слайд 3 — Актуальность
% ============================================================
\begin{frame}{Актуальность}

\begin{itemize}
  \item Растущая потребность в инструментах планирования задач
  \item Существующие решения~--- проприетарные, закрытый код
  \item Статистика и аналитика~--- только в платных версиях
  \item Нет возможности самостоятельного развёртывания
\end{itemize}

\vspace{0.5cm}
\textbf{Решение:} веб-приложение с открытым исходным кодом, REST~API, Swagger-документацией и встроенной статистикой.

\end{frame}

% ============================================================
% Слайд 4 — Сравнение аналогов
% ============================================================
\begin{frame}{Сравнение аналогов}

\centering
\footnotesize
\begin{tabular}{|l|c|c|c|c|c|}
\hline
\textbf{Критерий} & \textbf{Todoist} & \textbf{MS To~Do} & \textbf{TickTick} & \textbf{G.~Tasks} & \textbf{TodoApp} \\
\hline
Open Source & Нет & Нет & Нет & Нет & \textbf{Да} \\
\hline
REST API & Да & Да & Да & Да & \textbf{Да} \\
\hline
Статистика & Платно & Нет & Платно & Нет & \textbf{Да} \\
\hline
Swagger & Нет & Нет & Нет & Нет & \textbf{Да} \\
\hline
Self-hosted & Нет & Нет & Нет & Нет & \textbf{Да} \\
\hline
\end{tabular}

\end{frame}

% ============================================================
% Слайд 5 — Стек технологий
% ============================================================
\begin{frame}{Стек технологий}

\begin{columns}[T]
\begin{column}{0.5\textwidth}
\textbf{Основные:}
\begin{itemize}
  \item Go 1.24
  \item PostgreSQL 18
  \item Docker / Docker Compose
  \item net/http (стандартная библиотека)
\end{itemize}
\end{column}
\begin{column}{0.5\textwidth}
\textbf{Библиотеки:}
\begin{itemize}
  \item pgx 5.8 --- драйвер PostgreSQL
  \item zap --- логирование
  \item swaggo/swag --- Swagger
  \item validator --- валидация
  \item golang-migrate --- миграции
\end{itemize}
\end{column}
\end{columns}

\end{frame}

% ============================================================
% Слайд 6 — Архитектура
% ============================================================
\begin{frame}{Многослойная архитектура}

\centering
\begin{tikzpicture}[
  box/.style={rectangle, draw, fill=blue!10, minimum width=5cm, minimum height=0.9cm, text centered, font=\small},
  arr/.style={-{Stealth[length=3mm]}, thick}
]
\node[box] (transport) {Transport (HTTP handlers, middleware)};
\node[box, below=0.7cm of transport] (service) {Service (бизнес-логика)};
\node[box, below=0.7cm of service] (repository) {Repository (SQL-запросы)};
\node[box, below=0.7cm of repository, fill=green!10] (database) {PostgreSQL};

\draw[arr] (transport) -- (service);
\draw[arr] (service) -- (repository);
\draw[arr] (repository) -- (database);

\node[right=1cm of service, font=\footnotesize, text width=4cm] {Интерфейс репозитория определён в сервисе\\(Dependency Inversion)};
\end{tikzpicture}

\end{frame}

% ============================================================
% Слайд 7 — Модульная структура
% ============================================================
\begin{frame}{Модульная структура}

\centering
\begin{tikzpicture}[
  box/.style={rectangle, draw, fill=blue!10, minimum width=2.2cm, minimum height=0.7cm, text centered, font=\footnotesize},
  core/.style={rectangle, draw, fill=orange!15, minimum width=2.5cm, minimum height=0.7cm, text centered, font=\footnotesize},
  arr/.style={-{Stealth[length=2mm]}, thick}
]
\node[box] (main) {main.go};
\node[box, below left=1cm and 1cm of main] (users) {users};
\node[box, below left=1cm and -0.5cm of main] (tasks) {tasks};
\node[box, below right=1cm and -0.5cm of main] (stats) {statistics};
\node[box, below right=1cm and 1cm of main] (web) {web};
\node[core, below=2.5cm of main] (core) {core};

\draw[arr] (main) -- (users);
\draw[arr] (main) -- (tasks);
\draw[arr] (main) -- (stats);
\draw[arr] (main) -- (web);
\draw[arr] (users) -- (core);
\draw[arr] (tasks) -- (core);
\draw[arr] (stats) -- (core);
\draw[arr] (web) -- (core);
\end{tikzpicture}

\vspace{0.3cm}
Каждый модуль: transport $\rightarrow$ service $\rightarrow$ repository

\end{frame}

% ============================================================
% Слайд 8 — ER-диаграмма
% ============================================================
\begin{frame}{Проектирование базы данных}

\centering
\begin{tikzpicture}[
  entity/.style={rectangle, draw, minimum width=3.5cm, minimum height=0.5cm, font=\footnotesize\ttfamily},
  title/.style={rectangle, draw, fill=blue!15, minimum width=3.5cm, minimum height=0.5cm, font=\footnotesize\bfseries\ttfamily}
]
\node[title] (ut) {users};
\node[entity, below=0pt of ut] (u1) {id : UUID (PK)};
\node[entity, below=0pt of u1] (u2) {version : INT};
\node[entity, below=0pt of u2] (u3) {full\_name : VARCHAR};
\node[entity, below=0pt of u3] (u4) {phone\_number : VARCHAR};

\node[title, right=2.5cm of ut] (tt) {tasks};
\node[entity, below=0pt of tt] (t1) {id : UUID (PK)};
\node[entity, below=0pt of t1] (t2) {version : INT};
\node[entity, below=0pt of t2] (t3) {title : VARCHAR};
\node[entity, below=0pt of t3] (t4) {description : VARCHAR};
\node[entity, below=0pt of t4] (t5) {completed : BOOLEAN};
\node[entity, below=0pt of t5] (t6) {created\_at : TIMESTAMP};
\node[entity, below=0pt of t6] (t7) {completed\_at : TIMESTAMP};
\node[entity, below=0pt of t7] (t8) {author\_user\_id : UUID (FK)};

\draw[thick] (u1.east) -- ++(0.8cm,0) |- (t8.west) node[midway, above, font=\footnotesize] {1..*};
\end{tikzpicture}

\end{frame}

% ============================================================
% Слайд 9 — REST API
% ============================================================
\begin{frame}{REST API (11 эндпоинтов)}

\centering
\footnotesize
\begin{tabular}{|c|l|l|}
\hline
\textbf{Метод} & \textbf{Путь} & \textbf{Описание} \\
\hline
GET & /api/v1/users & Список пользователей \\
POST & /api/v1/users & Создать пользователя \\
GET & /api/v1/users/\{id\} & Получить пользователя \\
PATCH & /api/v1/users/\{id\} & Обновить пользователя \\
DELETE & /api/v1/users/\{id\} & Удалить пользователя \\
\hline
GET & /api/v1/tasks & Список задач \\
POST & /api/v1/tasks & Создать задачу \\
GET & /api/v1/tasks/\{id\} & Получить задачу \\
PATCH & /api/v1/tasks/\{id\} & Обновить задачу \\
DELETE & /api/v1/tasks/\{id\} & Удалить задачу \\
\hline
GET & /api/v1/statistics & Статистика \\
\hline
\end{tabular}

\end{frame}

% ============================================================
% Слайд 10 — Паттерны проектирования
% ============================================================
\begin{frame}{Паттерны проектирования}

\begin{description}[leftmargin=0.5cm]
  \item[Dependency Injection] --- зависимости передаются извне, ручная сборка в main.go
  \item[Repository Pattern] --- абстракция доступа к данным, интерфейс в пакете сервиса
  \item[Optimistic Locking] --- поле version, UPDATE \ldots{} WHERE version~=~\$2
  \item[Nullable{[T]}] --- различение <<не передано>>, <<null>>, <<значение>> в PATCH
\end{description}

\end{frame}

% ============================================================
% Слайд 11 — Пример кода
% ============================================================
\begin{frame}[fragile]{Пример: доменная сущность Task}

\begin{lstlisting}[language=Go]
type Task struct {
    ID           uuid.UUID
    Version      int
    Title        string
    Description  *string
    Completed    bool
    CreatedAt    time.Time
    CompletedAt  *time.Time
    AuthorUserID uuid.UUID
}

func CreateTask(title string, desc *string,
    authorID uuid.UUID) Task {
    return NewTask(uuid.New(), 1, title, desc,
        false, time.Now(), nil, authorID)
}
\end{lstlisting}

\end{frame}

% ============================================================
% Слайд 12 — Middleware
% ============================================================
\begin{frame}{Цепочка Middleware}

\centering
\begin{tikzpicture}[
  box/.style={rectangle, draw, fill=blue!8, minimum width=3cm, minimum height=0.6cm, text centered, font=\footnotesize},
  arr/.style={-{Stealth[length=2mm]}, thick}
]
\node[box] (cors) {CORS};
\node[box, below=0.3cm of cors] (rid) {RequestID};
\node[box, below=0.3cm of rid] (log) {Logger};
\node[box, below=0.3cm of log] (trace) {Trace};
\node[box, below=0.3cm of trace] (panic) {Panic Recovery};
\node[box, below=0.3cm of panic, fill=green!10] (handler) {Handler};

\draw[arr] (cors) -- (rid);
\draw[arr] (rid) -- (log);
\draw[arr] (log) -- (trace);
\draw[arr] (trace) -- (panic);
\draw[arr] (panic) -- (handler);
\end{tikzpicture}

\vspace{0.3cm}
Каждый middleware оборачивает следующий обработчик

\end{frame}

% ============================================================
% Слайд 13 — Тестирование
% ============================================================
\begin{frame}{Результаты тестирования}

\centering
\footnotesize
\begin{tabular}{|c|l|c|}
\hline
\textbf{№} & \textbf{Тестовый сценарий} & \textbf{Статус} \\
\hline
1 & Создание пользователя (валидные данные) & Пройден \\
2 & Создание пользователя (невалидный телефон) & Пройден \\
3 & Получение списка пользователей & Пройден \\
4 & Создание задачи & Пройден \\
5 & Задача без автора & Пройден \\
6 & Обновление статуса (completed = true) & Пройден \\
7 & Удаление задачи & Пройден \\
8 & Несуществующая задача (404) & Пройден \\
9 & Получение статистики & Пройден \\
10 & Конфликт версии (409) & Пройден \\
\hline
\end{tabular}

\vspace{0.3cm}
\textbf{10/10 тестов пройдены успешно}

\end{frame}

% ============================================================
% Слайд 14 — Скриншоты
% ============================================================
\begin{frame}{Демонстрация}

\begin{columns}[T]
\begin{column}{0.5\textwidth}
\centering
\fbox{\parbox{4.5cm}{\centering\vspace{2cm}\footnotesize\textit{Скриншот\\Swagger UI}\vspace{2cm}}}

Swagger UI
\end{column}
\begin{column}{0.5\textwidth}
\centering
\fbox{\parbox{4.5cm}{\centering\vspace{2cm}\footnotesize\textit{Скриншот\\веб-интерфейса}\vspace{2cm}}}

Веб-интерфейс TodoApp
\end{column}
\end{columns}

\end{frame}

% ============================================================
% Слайд 15 — Заключение
% ============================================================
\begin{frame}{Заключение}

\textbf{Реализовано:}
\begin{itemize}
  \item REST API на Go с многослойной архитектурой (11 эндпоинтов)
  \item CRUD пользователей и задач с валидацией
  \item Оптимистичная блокировка, graceful shutdown
  \item Статистика с фильтрацией
  \item Swagger-документация, веб-интерфейс
  \item Контейнеризация Docker Compose
\end{itemize}

\vspace{0.5cm}
\textbf{Технологии:} Go, PostgreSQL, Docker, pgx, zap, swaggo

\vspace{0.3cm}
\textbf{Код:} \url{https://github.com/milas1221/golang-todoapp}

\end{frame}

% ============================================================
% Слайд 16 — Спасибо за внимание
% ============================================================
\begin{frame}
\centering
\vfill
{\Huge Спасибо за внимание!}
\vfill
\end{frame}

\end{document}
